% Options for packages loaded elsewhere
\PassOptionsToPackage{unicode}{hyperref}
\PassOptionsToPackage{hyphens}{url}
\documentclass[
  a4paper,
]{ltjsarticle}
\usepackage{xcolor}
\usepackage[margin=25mm]{geometry}
\usepackage{amsmath,amssymb}
\setcounter{secnumdepth}{-\maxdimen} % remove section numbering
\usepackage{iftex}
\ifPDFTeX
  \usepackage[T1]{fontenc}
  \usepackage[utf8]{inputenc}
  \usepackage{textcomp} % provide euro and other symbols
\else % if luatex or xetex
  \usepackage{unicode-math} % this also loads fontspec
  \defaultfontfeatures{Scale=MatchLowercase}
  \defaultfontfeatures[\rmfamily]{Ligatures=TeX,Scale=1}
\fi
\usepackage{lmodern}
\ifPDFTeX\else
  % xetex/luatex font selection
\fi
% Use upquote if available, for straight quotes in verbatim environments
\IfFileExists{upquote.sty}{\usepackage{upquote}}{}
\IfFileExists{microtype.sty}{% use microtype if available
  \usepackage[]{microtype}
  \UseMicrotypeSet[protrusion]{basicmath} % disable protrusion for tt fonts
}{}
\makeatletter
\@ifundefined{KOMAClassName}{% if non-KOMA class
  \IfFileExists{parskip.sty}{%
    \usepackage{parskip}
  }{% else
    \setlength{\parindent}{0pt}
    \setlength{\parskip}{6pt plus 2pt minus 1pt}}
}{% if KOMA class
  \KOMAoptions{parskip=half}}
\makeatother
\usepackage{longtable,booktabs,array}
\newcounter{none} % for unnumbered tables
\usepackage{calc} % for calculating minipage widths
% Correct order of tables after \paragraph or \subparagraph
\usepackage{etoolbox}
\makeatletter
\patchcmd\longtable{\par}{\if@noskipsec\mbox{}\fi\par}{}{}
\makeatother
% Allow footnotes in longtable head/foot
\IfFileExists{footnotehyper.sty}{\usepackage{footnotehyper}}{\usepackage{footnote}}
\makesavenoteenv{longtable}
\usepackage{graphicx}
\makeatletter
\newsavebox\pandoc@box
\newcommand*\pandocbounded[1]{% scales image to fit in text height/width
  \sbox\pandoc@box{#1}%
  \Gscale@div\@tempa{\textheight}{\dimexpr\ht\pandoc@box+\dp\pandoc@box\relax}%
  \Gscale@div\@tempb{\linewidth}{\wd\pandoc@box}%
  \ifdim\@tempb\p@<\@tempa\p@\let\@tempa\@tempb\fi% select the smaller of both
  \ifdim\@tempa\p@<\p@\scalebox{\@tempa}{\usebox\pandoc@box}%
  \else\usebox{\pandoc@box}%
  \fi%
}
% Set default figure placement to htbp
\def\fps@figure{htbp}
\makeatother
\setlength{\emergencystretch}{3em} % prevent overfull lines
\providecommand{\tightlist}{%
  \setlength{\itemsep}{0pt}\setlength{\parskip}{0pt}}
% Glyph map for lualatex (newunicodechar). Maps symbols that may lack font glyphs.
\usepackage{newunicodechar}
\usepackage{amssymb}
\newunicodechar{∎}{\ensuremath{\blacksquare}}
\newunicodechar{✓}{\ensuremath{\checkmark}}
\newunicodechar{✗}{\ensuremath{\times}}
\newunicodechar{⟹}{\ensuremath{\Longrightarrow}}
\newunicodechar{⟺}{\ensuremath{\Longleftrightarrow}}
\newunicodechar{↪}{\ensuremath{\hookrightarrow}}
\newunicodechar{⌊}{\ensuremath{\lfloor}}
\newunicodechar{⌋}{\ensuremath{\rfloor}}
\newunicodechar{≃}{\ensuremath{\simeq}}
\newunicodechar{≈}{\ensuremath{\approx}}
\newunicodechar{≤}{\ensuremath{\leq}}
\newunicodechar{≥}{\ensuremath{\geq}}
\newunicodechar{≠}{\ensuremath{\neq}}
\newunicodechar{→}{\ensuremath{\rightarrow}}
\newunicodechar{←}{\ensuremath{\leftarrow}}
\newunicodechar{↔}{\ensuremath{\leftrightarrow}}
\newunicodechar{⊥}{\ensuremath{\perp}}
\newunicodechar{√}{\ensuremath{\surd}}
\newunicodechar{∑}{\ensuremath{\sum}}
\newunicodechar{∏}{\ensuremath{\prod}}
\newunicodechar{∞}{\ensuremath{\infty}}
\newunicodechar{∈}{\ensuremath{\in}}
\newunicodechar{∀}{\ensuremath{\forall}}
\newunicodechar{∣}{\ensuremath{\mid}}
\newunicodechar{γ}{\ensuremath{\gamma}}
\newunicodechar{λ}{\ensuremath{\lambda}}
\newunicodechar{μ}{\ensuremath{\mu}}
\newunicodechar{ρ}{\ensuremath{\rho}}
\newunicodechar{σ}{\ensuremath{\sigma}}
\newunicodechar{φ}{\ensuremath{\varphi}}
\newunicodechar{ψ}{\ensuremath{\psi}}
\newunicodechar{θ}{\ensuremath{\theta}}
\newunicodechar{Δ}{\ensuremath{\Delta}}
\newunicodechar{Π}{\ensuremath{\Pi}}
\newunicodechar{Λ}{\ensuremath{\Lambda}}
\newunicodechar{τ}{\ensuremath{\tau}}
\newunicodechar{ε}{\ensuremath{\varepsilon}}
\newunicodechar{ω}{\ensuremath{\omega}}
\newunicodechar{π}{\ensuremath{\pi}}
\newunicodechar{ν}{\ensuremath{\nu}}
\newunicodechar{±}{\ensuremath{\pm}}
\newunicodechar{×}{\ensuremath{\times}}
\newunicodechar{·}{\ensuremath{\cdot}}
\usepackage{bookmark}
\IfFileExists{xurl.sty}{\usepackage{xurl}}{} % add URL line breaks if available
\urlstyle{same}
\hypersetup{
  pdftitle={Post-Lock Geometry --- The Initial Plane and the Terminal Plane Span 4 Dimensions (Two Nonzero Principal Angles); Plane Switching Is Not an Aggregation Artifact of a Long-Window SVD},
  pdfauthor={Noriaki Kihara (WF System Co., Ltd.), ORCID 0009-0004-6753-4020},
  hidelinks,
  pdfcreator={LaTeX via pandoc}}

\title{Post-Lock Geometry --- The Initial Plane and the Terminal Plane
Span 4 Dimensions (Two Nonzero Principal Angles); Plane Switching Is Not
an Aggregation Artifact of a Long-Window SVD}
\author{Noriaki Kihara (WF System Co., Ltd.), ORCID 0009-0004-6753-4020}
\date{2026-09-12}

\begin{document}
\maketitle

\textbf{Series:} Foundations of Self-Consistent Relational-Wave Closed
Systems (Paper 7 of 10)\\
\textbf{Author:} Noriaki Kihara (WF System Co., Ltd.)　\textbf{Date:}
2026-09-12\\
\textbf{Version DOI:} 10.5281/zenodo.22729108\\
\textbf{Concept DOI:} 10.5281/zenodo.22729107\\
\textbf{ORCID:} 0009-0004-6753-4020\\
\textbf{Zenodo:} https://zenodo.org/records/22729108\\

\begin{center}\rule{0.5\linewidth}{0.5pt}\end{center}

\subsection{Abstract}\label{abstract}

We fix geometrically what inflation (Paper 6) and SSB lock (Paper 5) do
in state space. The map is a real orthogonal rotation \(z(t)=Q_t z_0\),
and the state remains on a 2-dimensional plane (rank 2) at every instant
(the singular values of \([\Re z,\Im z]\) take only two values at every
step, and \(z^{\mathsf T}z=0\) conservation gives an orthogonal
isometric frame). On the floor the floor plane
\(P_0=\mathrm{span}\{\Re z_0,\Im z_0\}\) is \(K\)-invariant and the
state stays in \(P_0\). When SSB leaves the floor, the lock plane
\(P_{\rm term}\) tilts away from \(P_0\), and the principal angles
between the two planes are, independently of \(N\), \textbf{both
nonzero}
(\(N=6{:}27.1^\circ,30.8^\circ\)／\(N=40{:}18.1^\circ,20.4^\circ\)), so
the dimension spanned by the two planes is \(4\). \(P_0\) and
\(P_{\rm term}\) span 4 dimensions with two nonzero principal angles,
and from the same floor the terminal orientation is seed-dependent and
not uniquely determined. We do not claim a specific \(SO(4)\) rotation
structure for the cumulative map (in the real data \(M_t\) is confirmed
not to preserve a 4-dimensional subspace). The tilt of the lock plane
away from the floor is, over all \(N=6\)--\(40\),
\(14.2^\circ\)--\(29.0^\circ\) (median \(\sim19.2^\circ\)). During the
transition a second direction rises (transient
\(s_3/s_1\approx0.10\)--\(0.20\) = effective rank 4), and after lock the
system re-locks to a single plane (tail
\(s_3/s_1\approx10^{-5}\)--\(10^{-13}\) = rank 2). By a short-window SVD
test we separate this ``plane switching'' from an appearance of
long-window aggregation, showing it is a \textbf{real plane tilt}.
Classification: derived geometry + measurement (artifact excluded).

\begin{center}\rule{0.5\linewidth}{0.5pt}\end{center}

\subsection{1. Problem}\label{problem}

Reading the endpoint merely as a ``circle'' fails to capture the
essential observation of this map --- \textbf{switching from one plane
to another plane}. We quantify this for all \(N\ge6\) and separate the
appearance (aggregation of a long-window SVD) from the real rotation.

\subsection{2. Instantaneous Structure --- Rank 2 (a Single
Plane)}\label{instantaneous-structure-rank-2-a-single-plane}

The map is a real orthogonal rotation \(z(t)=Q_t z_0\) (\(Q_t\) real
orthogonal, acting identically on \(\Re z,\Im z\)), and the state stays
on a 2-dimensional plane at every instant (the singular values of
\([\Re z,\Im z]\in\mathbb R^{2\times M}\) take only two values at every
step, and \(z^{\mathsf T}z=0\) conservation gives an orthogonal
isometric frame). \textbf{Formulation as Grassmann motion}: for
\(z=a+ib\), from \(z^{\mathsf T}z=0\) and \(\|z\|=1\) we have
\(\|a\|^2=\|b\|^2=\tfrac12\) and \(a^{\mathsf T}b=0\), so
\(u=\sqrt2\,a,\ v=\sqrt2\,b\) are always an orthonormal real 2-frame
\((u,v)\in V_2(\mathbb R^M)\). The global phase
\(z\mapsto e^{i\theta}z\) merely rotates \((u,v)\) within the same
plane, so the object observed --- modulo global phase --- is the
\textbf{2-dimensional subspace \(P(t)=\mathrm{span}\{u,v\}\) itself}
(motion on the Grassmann manifold \(\mathrm{Gr}(2,M)\)). The projection
\(\Pi_t=u_tu_t^{\mathsf T}+v_tv_t^{\mathsf T}\) moves exactly, by each
step's real orthogonal \(Q_t=e^{\Delta\tau K_t}\), as

\[\Pi_{t+1}=Q_t\,\Pi_t\,Q_t^{\mathsf T}\]

On the floor \(Q_0P_0=P_0\), so \(\Pi_{t+1}=\Pi_t\) (relative
equilibrium, fixed); in inflation \(\Pi_t\neq\Pi_0\) and the plane
moves; at saturation it re-locks to \(P_{\rm term}\). Sections 3--5
measure this plane motion \(\Pi_0\to\Pi_t\to\Pi_{\rm term}\). On the
floor \(P_0=\mathrm{span}\{\Re z_0,\Im z_0\}\) is \(K\)-invariant
(\(Ka=\mu b,\ Kb=-\mu a\)) and the state stays in \(P_0\) (principal
angles zero). The state-space SVD ratio \(s_3/s_1\) of the third/first
singular values over the last 200 steps after lock is
\(10^{-5}\)--\(10^{-13}\) = a single plane (rank 2/\(S^1\)).

\subsection{3. The Dimension Spanned by the Initial Plane and the
Terminal Plane --- 4D
Span}\label{the-dimension-spanned-by-the-initial-plane-and-the-terminal-plane-4d-span}

When SSB leaves the floor, the lock plane \(P_{\rm term}\) tilts away
from \(P_0\). The principal angles between the initial plane \(P_0\) and
the terminal plane \(P_{\rm term}\) are, independently of \(N\),
\textbf{both nonzero}
(\(N=6{:}27.1^\circ,30.8^\circ\)／\(N=40{:}18.1^\circ,20.4^\circ\)), so
the dimension spanned by the two planes is \(4\). The tilt of the lock
plane away from the floor (\(\arcsin\sqrt{H_\perp/H_{\rm end}}\)) is,
over all \(N=6\)--\(40\), as follows (excerpted from
\texttt{plane\_switch\_allN.csv}):

{\def\LTcaptype{none} % do not increment counter
\begin{longtable}[]{@{}llll@{}}
\toprule\noalign{}
\(N\) & tilt {[}°{]} & transient \(s_3/s_1\) & tail \(s_3/s_1\) \\
\midrule\noalign{}
\endhead
\bottomrule\noalign{}
\endlastfoot
6 & 29.0 & 0.172 & \(3.8\times10^{-6}\) \\
8 & 20.3 & 0.135 & \(2.8\times10^{-4}\) \\
16 & 24.5 & 0.172 & \(1.1\times10^{-5}\) \\
22 & 15.7 & 0.130 & \(1.1\times10^{-13}\) \\
40 & 19.3 & 0.132 & \(4.5\times10^{-7}\) \\
64 & 21.7 & 0.021 & \(2.5\times10^{-4}\) \\
\end{longtable}
}

Over all \(N\) the tilt is \(14.2^\circ\)--\(29.0^\circ\) (median
\(\sim19.2^\circ\)), neither \(0^\circ\) nor \(90^\circ\) = isotropy is
broken.

\textbf{The tilt, the two principal angles, and \(H_\perp/H\) are not
independent measured quantities but the same Grassmann geometry.} For
the initial and terminal frames \(U_0=(u_0,v_0),\ U_1=(u_1,v_1)\) and
\(P_0^\perp=I-U_0U_0^{\mathsf T}\), from the definition of principal
angles (the singular values of \(U_0^{\mathsf T}U_1\) equal
\(\cos\theta_i\)) we have
\(\|P_0^\perp U_1\|_F^2=\sin^2\theta_1+\sin^2\theta_2\). Since
\(z=\tfrac1{\sqrt2}(u+iv)\),

\[\frac{H_\perp}{H}=\frac12\left(\sin^2\theta_1+\sin^2\theta_2\right)=\frac14\big\|\Pi_{\rm term}-\Pi_0\big\|_F^2,\qquad \alpha=\arcsin\sqrt{\frac{\sin^2\theta_1+\sin^2\theta_2}{2}}\]

(from \(\Pi=UU^{\mathsf T}\) we get
\(\|\Pi_1-\Pi_0\|_F^2=2(\sin^2\theta_1+\sin^2\theta_2)\)). This is not
an additional hypothesis but an \textbf{exact geometric identity} from
the definitions, and \(N=6\)'s
\((27.1^\circ,30.8^\circ)\to\alpha=29.0^\circ\) and \(N=40\)'s
\((18.1^\circ,20.4^\circ)\to\alpha=19.3^\circ\) agree with the table.
That is, the tilt \(\alpha\) away from the floor is the
\textbf{Grassmann distance} \(\tfrac14\|\Pi_{\rm term}-\Pi_0\|_F^2\)
from the initial plane itself.

\begin{figure}
\centering
\pandocbounded{\includegraphics[keepaspectratio,alt={Figure 1: Plane switching (all N≥6). Left = tilt of the lock plane away from the floor vs N; right = transient s3/s1 (rank 4/S³) → tail s3/s1 (rank 2/S¹).}]{figs/p7_postlock_geometry_en_fig1.png}}
\caption{Figure 1: Plane switching (all N≥6). Left = tilt of the lock
plane away from the floor vs N; right = transient s3/s1 (rank 4/S³) →
tail s3/s1 (rank 2/S¹).}
\end{figure}

\textbf{Figure 1} \(N\)-dependence of plane switching.

\subsection{4. The Initial Plane and the Terminal Plane Span 4
Dimensions}\label{the-initial-plane-and-the-terminal-plane-span-4-dimensions}

\(P_0\) and \(P_{\rm term}\) meet at two nonzero principal angles
(\(N=6{:}27.1^\circ,30.8^\circ\)／\(N=40{:}18.1^\circ,20.4^\circ\)) and
\(\dim(P_0+P_{\rm term})=4\) --- the terminal plane tilts away from the
initial plane in two independent directions (these ``4 directions'' are
meant only in the sense of the dimension of the span). \textbf{From the
same floor this terminal orientation is seed-dependent and not uniquely
determined} (Paper 5). We do not claim a specific \(SO(4)\) rotation
structure for the cumulative map
\(M_t=\prod_n\exp(\Delta\tau K(\varphi_n))\) --- in the real data
\(M_t\) is confirmed not to preserve this 4-dimensional subspace (the
invariance \(M_tS=S\) is broken), and \(M_t\) is an orthogonal rotation
of the full \(\mathbb R^M\) that tilts \(P_0\) into the 4-dimensional
span.

During the transition a second rotation plane rises (transient
\(s_3/s_1\approx0.10\)--\(0.20\) = rank 4), but after lock the system
re-locks to a single plane (tail \(s_3/s_1\approx0\) = rank 2). At
\(N=64\) the transient \(s_3/s_1=0.021\) is small, so the second-plane
transient weakens for large \(N\). This transitional rank 4 is
temporary; after lock it returns to rank 2, and the cumulative map also
does not preserve a 4-dimensional subspace (§4) --- \textbf{no
persistent \(S^3\) structure is found} (inspected, short-window SVD +
discriminating computation).

In the time-varying, non-normal tangent cocycle confirmed in Paper 6,
the maximally amplified direction rotates with time while it is
amplified. The additional direction of the transient observed in this
paper can be read as the geometric manifestation of that tangent
dynamics tilting the instantaneous rank-2 plane \(P(t)\) with time.
Therefore the effective rank 4 of the transient is \textbf{not a
4-dimensionalization of the instantaneous state but the local span swept
in a short time by the moving 2-plane \(P(t)\)}, and after saturation it
re-locks at rank 2 to a single terminal plane. That is, Paper 6 (in
tangent space the amplified direction rotates) and Paper 7 (in state
space the 2-plane \(P(t)\) tilts and moves to another 2-plane) are the
tangent-side and state-side representations of the same phenomenon.
Using \(H_\perp/H=\tfrac14\|\Pi_t-\Pi_0\|_F^2\) (the §3 identity), Paper
6's initial ramp \(H_\perp/H\propto e^{2\gamma t}\) is geometrically
\(\|\Pi_t-\Pi_0\|_F\propto e^{\gamma t}\) --- \textbf{inflation is an
exponential departure from the initial plane on the Grassmann space
\(\mathrm{Gr}(2,M)\)}.

\subsection{5. Exclusion of the
Artifact}\label{exclusion-of-the-artifact}

A long-window SVD can, if a single plane rotates with time, be
aggregated into an apparently higher rank. With the short-window SVD
test (\texttt{artifact\_test\_shortwindow\_svd}) we separated the fact
that, instantaneously staying at rank 2, the initial plane \(P_0\) tilts
in two directions (\(P_0\) and \(P_{\rm term}\) span 4 dimensions) ---
the ``plane switching'' is not an appearance of long-window aggregation
but a \textbf{real plane tilt}.

\begin{figure}
\centering
\pandocbounded{\includegraphics[keepaspectratio,alt={Figure 2: Short-window SVD test. While preserving instantaneous rank 2, the initial plane tilts in two directions (plane switching is real).}]{figs/p7_postlock_geometry_en_fig2.png}}
\caption{Figure 2: Short-window SVD test. While preserving instantaneous
rank 2, the initial plane tilts in two directions (plane switching is
real).}
\end{figure}

\textbf{Figure 2} Artifact exclusion (short-window SVD).

\subsection{6. A Caveat on Readout}\label{a-caveat-on-readout}

When there is initially only a single complex plane, its normal is a
degenerate orthogonal complement and is unobservable (conceptual). Only
once inflation raises a specific direction in the normal space does the
relative phase (tilt) between the initial plane and the new plane become
readable. \textbf{What can be read out is always a relational quantity
(the relative phase between planes)}; the absolute normal direction
cannot be recovered from the state alone (\(Q_t\) hides the orientation,
and there is no conserved quantity pointing to it).

\subsection{7. Claim Classification and Open
Items}\label{claim-classification-and-open-items}

\begin{itemize}
\tightlist
\item
  Classification: derived geometry (real orthogonal rotation rank 2) +
  measurement (principal angles, tilt, 4D span, artifact exclusion) +
  discriminating computation (confirmed that the cumulative map does not
  preserve a 4-dimensional subspace and is not \(SO(4)\)). Verdict: the
  4D span, the two nonzero principal angles, and the double rotation of
  the relative configuration are retained; ``the dynamics being an
  \(SO(4)\) double rotation on 4 dimensions'' is rejected.
\item
  Reference (a different regime): the crystalline terminal states at
  \(N=3,4\) are special solutions in which the plane refreezes at the
  magic angle \(\arccos\sqrt{2/3}=35.264^\circ\) (the general \(N\ge6\)
  lock is \(14\)--\(29^\circ\)). The crystal/glass bifurcation is a
  separate topic of this series.
\item
  Open: the physical identification of the 4D span (\(xyzt\) or
  \(xyztRQ\), etc.) and the type of gauge (scale) are outside the scope
  of this paper (Paper 0 §Open).
\end{itemize}

\subsection{Related Work (Structural Agreement, Not the Source of
Derivation)}\label{related-work-structural-agreement-not-the-source-of-derivation}

\begin{itemize}
\tightlist
\item
  The double rotation of \(SO(4)\) (two invariant planes, two
  independent rotation angles): this is a description of the
  \textbf{relative configuration} of \(P_0\to P_{\rm term}\), not a
  property of the cumulative map itself (in the §4 discriminating
  computation \(M_tS\neq S\)).
\item
  Relative equilibrium (Marsden \& Ratiu): the rigid-body-rotation
  relative equilibrium of the floor and the lock.
\end{itemize}

\subsection{References}\label{references}

\begin{enumerate}
\def\labelenumi{\arabic{enumi}.}
\tightlist
\item
  Noriaki Kihara, ``The mechanism of inflationary rapid expansion in
  self-consistent relational-wave closed systems,'' Concept DOI
  10.5281/zenodo.22112008.
\item
  J. E. Marsden and T. S. Ratiu, \emph{Introduction to Mechanics and
  Symmetry}, Springer (1999).
\end{enumerate}

\subsection{Reproduction}\label{reproduction}

The analysis scripts, figures, and CSV are in
\texttt{位相ロック全N確認\_相対平衡\_20260912/面の乗り換え全N\_20260912/}
(\texttt{analyze\_plane\_switch\_allN\_20260912.py},
\texttt{artifact\_test\_shortwindow\_svd\_20260912.py},
\texttt{plane\_switch\_allN.csv}, \texttt{fig\_plane\_switch\_allN.png},
\texttt{fig\_artifact\_shortwindow\_svd.png}). For the locked data,
small \(N\) and \(N=23\)--\(29\) are 500 steps, and the unsaturated
\(N\) (22, 30--40) and \(N=64\) are 2000 steps. The principal angles of
the initial and terminal planes are the step0/step2000 complex planes in
\texttt{N22\_2000step検証走行\_/} and \texttt{N40\_2000step検証走行\_/}.
The early rank-3 audit is in
\texttt{independent\_rank3\_geometry\_audit\_20260831/}. The dynamics
map is identical to the canonical \texttt{run\_N3\_N40\_stage123\_v1.py}
(SHA256
\texttt{1abf2353fee2e4f56f05e7a6f149fd086885136beb61ab571b48a56b09691567}),
readout only.

\end{document}
